\section{Four data classes that must not be conflated}
\label{sec:fourclasses}

A representation label and one cubic incidence do not define a physical RG
trajectory.  The finite obstruction can be stated without a parameter scan.
For a radial spectral potential
\begin{equation}
 V(x)=a_2x^2+a_3x^3+a_4x^4,
 \label{eq:radialpotential}
\end{equation}
four logically different data are required.

\begin{enumerate}[label=(\roman*)]
\item A \emph{moment character} fixes ratios such as
  $a_3^2/(a_2a_4)$ and the relative weight of single-trace and tensor-square
  invariants.  A common overall action unit cannot fix these ratios.
\item \emph{Pointed occupation} excludes the zero section.  Every unpointed
  polynomial of the form \eqref{eq:radialpotential} has $V'(0)=0$; covariance
  and positivity alone therefore do not force a nonzero occupied branch.
\item A \emph{physical scale} maps one dimensionless parent gap to a measured
  mass unit.  Rescaling all parent gaps preserves the dimensionless packet.
\item A \emph{renormalization map} specifies fields, gauge fixing, ghosts,
  scheme, thresholds and matching.  Finite coupling redefinitions change the
  coordinate representation of a beta vector even when observables agree.
\end{enumerate}

This separation prevents four common overclaims: normalization is not a
moment selector; spontaneous occupation is not implied by representation
content; a dimensionless scale ratio is not an SI prediction; and an exact
functional equation is not an evaluated component flow.

\begin{proposition}[Finite four-class obstruction]
The current scalar-15 edge, common unit and invariant support do not by
themselves determine a unique physical RG initial-value problem.
\end{proposition}
\begin{proof}
Multiplication by a common unit leaves $a_3^2/(a_2a_4)$ invariant.  Equation
\eqref{eq:radialpotential} has zero first derivative at the origin.  A common
rescaling of all dimensionful gaps leaves every normalized source relation
unchanged.  Finally, under $g'=g+cg^2$ a coordinate beta function transforms
as $\beta_{g'}=(1+2cg)\beta_g$.  These four independent transformations leave
the previously reconstructed packet intact while changing respectively its
moment shape, occupation verdict, SI scale or flow coordinates.
\end{proof}

\section{Conditional common origin of the four classes}
\label{sec:commonorigin}

The four classes need not be four unrelated additions.  One conditional
source architecture can own all four, provided its transport to the
scalar-15 operator is proved.

First, an oriented primitive response character supplies the moment sequence
\begin{equation}
 e^x=1+x+\frac{x^2}{2}+\frac{x^3}{6}+\frac{x^4}{24}+\cdots .
 \label{eq:momentcharacter}
\end{equation}
Transport of this one character to $\mathcal D_{15}$ fixes relative spectral
moments; its two-copy multiplicative character supplies the tensor-square
sector.  Reversal averaging cannot be imposed before the odd cubic has
selected an oriented branch, because doing so would delete the very incidence
that induces the scalar response.

Second, a pointed seed orientation may activate an already selected Higgs
body coordinate,
\begin{equation}
 V_H(h)=-k_Hh^2+\lambda_Hh^4,
 \qquad h_*^2=\frac{k_H}{2\lambda_H}>0.
\end{equation}
The positive matching row of Eq.~\eqref{eq:matchingrow} then yields
\begin{equation}
 \boxed{\Delta_*=\frac{\mu_\Delta}{K_\Delta}h_*^2
 =\frac{k_H\mu_\Delta}{2K_\Delta\lambda_H}\ne0.}
 \label{eq:pointedoccupation}
\end{equation}
This is not an external gauge-charged linear source: the seed selects an
oriented orbit, the body occupies the Higgs coordinate, and the covariant
two-Higgs edge induces the scalar coordinate.

Third, write the physical operator as
\begin{equation}
 \mathcal D_{15}^{\rm phys}=M_*\widehat{\mathcal D}_{15}.
 \label{eq:physicalgap}
\end{equation}
This identifies where the common scale lives while retaining the metrological
boundary stated in Section~\ref{sec:scalar15handoff}.

Fourth, for observer resolution $k$ the covariant spectral family
\begin{equation}
 P_k(\mathcal D_{15})=\mathbf 1_{[-k,k]}(\mathcal D_{15}),
 \qquad
 \Gamma_k=\operatorname{Desc}_{P_k}\Gamma
 \label{eq:spectraldescent}
\end{equation}
is nested and gauge covariant:
$P_{k_1}P_{k_2}=P_{k_1}$ for $k_1\le k_2$, and
$P_k(W\mathcal DW^\dagger)=WP_k(\mathcal D)W^\dagger$.
Eigenvalue crossings therefore define source-owned threshold events rather
than thresholds inserted after a fit.  The smooth regulator used later is a
computational realization of the same resolution logic, not a new physical
coupling.

\begin{theorem}[Conditional four-origin assembly]
If the physical base--seed completion owns the oriented character transport,
the pointed Higgs cycle, one positive common gap and the covariant spectral
descent, then moment character, nonzero scalar occupation, relative scale and
renormalization prescription descend from one source architecture.
\end{theorem}
\begin{proof}
Equations \eqref{eq:momentcharacter}, \eqref{eq:pointedoccupation},
\eqref{eq:physicalgap} and \eqref{eq:spectraldescent} supply the four data
classes of Section~\ref{sec:fourclasses}.  Each construction is covariant on
its declared carrier, and none introduces a terminal-specific gain.  The
conclusion is conditional because occupation of this complete source packet
by the unrestricted physical parent is not derived.
\end{proof}

\section{Bosonic capacity is a viability test, not selection}

For active field-dependent modes $m_j^2=k_jr^2$, the one-loop logarithmic
weight is
\begin{equation}
 64\pi^2 B=\sum_j(-1)^{F_j}n_jk_j^2.
\end{equation}
At the frozen source point used by the audit, $g=1$ and
$y_u=0.9615522319$, the gauge--fermion deficit is
\begin{equation}
 D_{GF}=3y_u^4-\frac{171}{200}=1.70955948.
\end{equation}
A complex scalar $15_H$ has thirty real modes.  If all active modes share the
portal jet $\kappa_{H\Delta}$, their positive contribution is
\begin{equation}
 \Delta(64\pi^2B)_{15}=30\kappa_{H\Delta}^2,
\end{equation}
so one copy can reverse the sign whenever
\begin{equation}
 \boxed{|\kappa_{H\Delta}|>\sqrt{D_{GF}/30}=0.238716.}
 \label{eq:capacitybound}
\end{equation}
Equation~\eqref{eq:capacitybound} establishes raw source-scale capacity.  It
does not select the portal, prove a Coleman--Weinberg vacuum or remove
decoupling and threshold effects.

\section{Equal-upstream RG counterfamily}

The representation fixes the contribution of one active complex $15_H$ to
the GUT-normalized one-loop gauge coefficients,
\begin{equation}
 \boxed{\Delta b^{15_H}=\left(\frac76,\frac76,\frac76\right).}
 \label{eq:gaugeshift}
\end{equation}
Thus a degenerate complete multiplet changes common running but not relative
one-loop gauge slopes.  A physical trajectory additionally requires portals,
independent scalar quartics, Yukawa and cubic data, component thresholds and
breaking-scale matching.

A controlled equal-upstream audit makes the missing information visible.  At
one common active interval, two admissible portal choices give
\begin{center}
\begin{tabular}{@{}ccc@{}}
\toprule
$\kappa_{H\Delta}$ & $\min(\lambda+\beta_\lambda/4)$ & zero present?\\
\midrule
$0.80$ & $+0.071386$ & no\\
$1.20$ & $-0.660165$ & yes\\
\bottomrule
\end{tabular}
\end{center}
while all already frozen upstream data are held fixed.  This is a finite
non-identifiability witness: choosing either trajectory without an additional
source law would be a fit, not a Tau prediction.

\section{Exact finite stopping point}

The scalar-15 branch and its upstream renormalization interface close at five distinct levels once the 14 November 2026 three-pillar theorems (\texttt{U1}, \texttt{U2}, \texttt{F1}, \texttt{F3}, \texttt{HOMOG-K4}) are incorporated:
\begin{center}\small
\begin{tabularx}{0.96\linewidth}{@{}p{0.28\linewidth}p{0.34\linewidth}X@{}}
\toprule
Question & Proved result & Remaining boundary\\
\midrule
Source preimage & Unique in the declared recovery-complete type-II class &
Unrestricted physical scalar-15 branch occupation\\
Spectral/Clifford seed scale (\texttt{U1}, \texttt{F1}) & $\widehat{\Delta}_{K_4}=\frac{4}{3}P_\perp=G_{\mathrm{Cl}}$ unifies \texttt{C2} with \texttt{C1} at $r_{\mathrm{seed}}=1/4$; $\Pi^2=\Pi\implies\dot{r}\equiv 0$ &
Occupied non-triviality postulate (\texttt{POST-OCC-NU1})\\
Intensive UV ratios (\texttt{U2}) & $a_{\mathrm{UV}}^{-2}$ renormalizes $1/(4G_{\mathrm{ren}})$; $M_*/E_P=1/\sqrt{24\pi}$ and $\eta_{\mathrm{hor}}\equiv 1$ have \texttt{0.00e+00} UV drift &
One CODATA SI scale anchor\\
Simplicial \& 60-decade RG protection (\texttt{HOMOG-K4}, \texttt{F3}) & Local $S_4\times S_4\implies\kappa_{\mathrm{dir}}^{\mathrm{bare}}\equiv 0$; $\operatorname{Tr}(T_{\mathrm{def}})=0$ + Frullani tail universality &
Component scalar-15 beta functions and truncation error\\
Phenomenological trajectory & Exact covariant flow; counterfamily shows non-uniqueness from unit alone &
Portal, quartic, threshold and matching data\\
\bottomrule
\end{tabularx}
\end{center}

Crucially, the apparent upstream RG fragilities flagged in earlier audits (\texttt{F41, F44, F47}) are resolved without parameter tuning:
\begin{enumerate}
\item \emph{Intensive UV-invariance vs.\ extensive $1/(4G_{\mathrm{ren}})$ absorption (Theorem~\texttt{U2}):} While Srednicki entanglement entropy scales extensively as $c_S A_{\mathrm{AH}}/a_{\mathrm{UV}}^2$ under radial lattice refinement, the induced gravitational counterterm shifts $1/(4G_{\mathrm{ren}})=1/(4G_{\mathrm{bare}})+c_S/(\pi a_{\mathrm{UV}}^2)$ with the identical coefficient, leaving all intensive dimensionless ratios---including $M_*/E_P=1/\sqrt{24\pi}$, $\beta_{\mathrm{KMS}}|\kappa_{\mathrm{KH}}|/(2\pi c)=1$, and $\eta_{\mathrm{hor}}\equiv 1$---identically UV-invariant (\texttt{0.00e+00} drift).
\item \emph{Laplacian--Clifford identity, projector normalcy, and vanishing bare spectator stiffness (Theorems~\texttt{U1}, \texttt{F1}, \texttt{HOMOG-K4 H1--H2}):} Chung's normalized Laplacian on regular $K_4$, $\widehat{\Delta}_{K_4}=D^{-1/2}L_{K_4}D^{-1/2}=\frac{4}{3}P_\perp=G_{\mathrm{Cl}}$, eliminates the \texttt{C1/C2} norm bifurcation at $r_{\mathrm{seed}}=1/4$, while pure-state idempotence $\Pi^2=\Pi$ ($\dot{r}\equiv 0$) renders $\Gamma'_{1\text{-loop}}(1/4)<0$ normal to $\mathcal{P}_1(M_2(\mathbb{C}))$ alongside tangential $S^2$ Hessian $+4I_2\succ 0$. On periodic Freudenthal--Kuhn $K_4$ tilings, local $S_4\times S_4$ gauge invariance enforces $\kappa_{\mathrm{dir}}^{\mathrm{bare}}\equiv 0$ identically (\texttt{R4GATE-NATN1} closed), and defect trace-freeness $\operatorname{Tr}(T_{\mathrm{def}})=0$ protects the 60-decade flow (\texttt{F3}).
\end{enumerate}
This establishes a clean finite stopping point: neither normalizing the packet nor scanning RG portals can decide whether Nature occupies the scalar-15/type-II branch rather than an alternative neutrino/matter completion.
